\input{preamble.tex}

\title{\huge{Oblig 2}}
\author{\LARGE{Thobias Høivik}}
\date{}

\begin{document}
\maketitle

\newpage
\begin{prob*}[1]
    We will prove the following statement:

    \[ 
        \text{Every separable Hilbert space} \ H \ \text{has an 
        orthonormal basis.}
    \]

    \begin{enumerate}[label=\alph*)]
        \item Let $D = \{u_1,u_2,\ldots\}$ be a dense subset 
            of $H$ (why does such a set exist?). 
            Show that there is a subset $D' \subseteq D$ with 
            elements $D' = \{v_1,v_2,\ldots\}$, where 
            $(v_1,v_2,\ldots)$ is linearly independent and 
            $\operatorname{span}(v_1,v_2,\ldots)$ is dense in 
            $H$.
        \item Apply the Gram-Schmidt orthogonalization procedure 
            to the list $(v_1,v_2,\ldots)$. Explain why the resulting 
            list $(e_1,e_2,\ldots)$ is orthonormal and its span 
            is dense in $H$.
        \item Prove that $(e_1,e_2,\ldots)$ is an orthonormal 
            basis for $H$.
    \end{enumerate}
\end{prob*}
\begin{proof}[Proof of (a)]
    Trivially, every hilbert space $H$ is dense in itself since 
    $\overline H = H$. In our case we are given the assumption 
    that $H$ is separable, so it follows that there exists 
    a countable dense $D \subseteq H$.

    We will construct $D' \subseteq D$ with the desired properties. 
    Let $v_1$ be the first element of $D$ which is non-zero. 
    Suppose we have chosen $n$ elements $D_n' = \{v_1,\ldots,v_n\}$. 
    Let $D_{n+1}' = D_n' \cup \{u_i \in D\setminus D_n'\}$ where 
    $u_i$ is the first element of $D \setminus D_n'$ which is 
    not in $\operatorname{span}(D_n')$. Let 
    \[
        D' := \bigcup_{n \in \mathbb N} D'_{n}.
    \]
    Then $D'$ is linearly independent by construction.

    Now $D \subseteq \operatorname{span}(D')$ because if 
    $u_i \not\in D'$ then 
    $u_i \in \operatorname{span}(D'_n)$ for some $n$, so 
    $u_i \in \operatorname{span}(D')$. If $u_i \in D'$ the inclusion 
    is trivial. 
    Now we see that 
    \[
        H = \overline D \subseteq \overline{\operatorname{span}(D')} 
        \subseteq H,
    \]
    so $D'$ is dense in $H$.
\end{proof}

\begin{proof}[Proof of (b)]
    Apply the Gram-Schmidt procedure to the linearly independent list
    $(v_1,v_2,\ldots)$. Define
    \[
        w_1 = v_1, \qquad e_1 = \frac{w_1}{\|w_1\|},
    \]
    and recursively
    \[
        w_n = v_n - \sum_{k=1}^{n-1} \langle v_n,e_k\rangle e_k,
        \qquad
        e_n = \frac{w_n}{\|w_n\|}.
    \]
    Since $(v_1,v_2,\ldots)$ is linearly independent, we have
    $w_n \neq 0$ for every $n$, so each $e_n$ is well-defined.

    By the Gram-Schmidt construction, the vectors $e_1,e_2,\ldots$
    are orthonormal. Moreover, for every $n$,
    \[
        \operatorname{span}(e_1,\ldots,e_n)
        =
        \operatorname{span}(v_1,\ldots,v_n).
    \]
    Hence
    \[
        \operatorname{span}(e_1,e_2,\ldots)
        =
        \operatorname{span}(v_1,v_2,\ldots).
    \]
    But by part (a), $\operatorname{span}(v_1,v_2,\ldots)$ is dense
    in $H$. Therefore $\operatorname{span}(e_1,e_2,\ldots)$ is also
    dense in $H$.
\end{proof}

\begin{proof}[Proof of (c)]
    We proved orthonormality in (b), and furthermore we showed 
    $H = \operatorname{span}(D')$ in (a). 
    Since $\operatorname{span}(D') = 
    \operatorname{span}(e_1,e_2\ldots)$
    we conclude that $(e_1,e_2,\ldots)$.
\end{proof}



\end{document}
